\subsection{Joins}
Sometimes, we need to compare values across two tables, or across two rows of the same table. An inner join creates a table whose rows contain every possible combination of rows from the source tables. We can then filter those combinations to find relationships between rows.

An inner join is created by specifying multiple source tables in a \lstinline{FROM} clause. For example, \lstinline{SELECT * FROM team AS a, team AS b;} combines each row of \lstinline{team} with every row of \lstinline{team}. The aliases \lstinline{a} and \lstinline{b} distinguish the two copies.

To find pairs of people with the same birth month:
\begin{lstlisting}
SELECT a.name, b.name FROM team AS a, team AS b
    WHERE a.name < b.name AND a.birth_month = b.birth_month;
\end{lstlisting}
The condition \lstinline{a.name < b.name} lists each pair once.

\subsection{Aggregation}
Aggregation combines information from multiple rows into a result. Functions such as \lstinline{MAX}, \lstinline{MIN}, \lstinline{COUNT}, and \lstinline{SUM} can summarize table data. For example, this counts the rows in \lstinline{team}:
\begin{lstlisting}
SELECT COUNT(*) FROM team;
\end{lstlisting}

\subsubsection{\lstinline{GROUP BY}}
\lstinline{GROUP BY} groups together rows with the same value in a column. An aggregation is then performed for each group, producing one result row per group. This query counts teammates for each pet type:
\begin{lstlisting}
SELECT pet, COUNT(*) FROM team GROUP BY pet;
\end{lstlisting}

\subsubsection{\lstinline{HAVING}}
\lstinline{HAVING} filters groups, just as \lstinline{WHERE} filters individual rows. This selects pet types owned by more than one teammate:
\begin{lstlisting}
SELECT pet FROM team GROUP BY pet HAVING COUNT(*) > 1;
\end{lstlisting}
